% ============================================================
%  《高中化学系统讲义》导言区
%  编译：xelatex（两遍）  字体：Noto Serif/Sans CJK SC + TeX Gyre Termes
% ============================================================
\usepackage[a4paper,top=2.1cm,bottom=2.0cm,left=1.9cm,right=1.9cm,
  headheight=14pt,headsep=0.45cm,footskip=0.75cm]{geometry}

% ---------- 字体 ----------
\usepackage{amsmath,amssymb}
\usepackage{unicode-math}
\setmainfont{TeX Gyre Termes}
\setsansfont{TeX Gyre Heros}
\setmathfont{TeX Gyre Termes Math}
\setCJKmainfont{Noto Serif CJK SC}[BoldFont={Noto Sans CJK SC Bold},ItalicFont={AR PL KaitiM GB}]
\setCJKsansfont{Noto Sans CJK SC}[BoldFont={Noto Sans CJK SC Bold}]
\setCJKmonofont{Noto Sans Mono CJK SC}
\setCJKfamilyfont{zhsong}{Noto Serif CJK SC}[BoldFont={Noto Serif CJK SC Bold}]
\setCJKfamilyfont{zhhei}{Noto Sans CJK SC}[BoldFont={Noto Sans CJK SC Bold}]
\setCJKfamilyfont{zhkai}{AR PL KaitiM GB}[AutoFakeBold=2]
\setCJKfamilyfont{zhfs}{AR PL KaitiM GB}
\def\songti{\CJKfamily{zhsong}}
\def\heiti{\CJKfamily{zhhei}}
\def\kaishu{\CJKfamily{zhkai}}
\linespread{1.22}
\xeCJKsetcharclass{"2460}{"24FF}{1}   % ①②③ 等按中文字符处理
\xeCJKsetcharclass{"2160}{"217F}{1}   % Ⅰ Ⅱ Ⅲ 罗马数字
\xeCJKsetcharclass{"25A0}{"25FF}{1}   % △ ○ 等几何符号用中文字体
\xeCJKsetcharclass{"00B7}{"00B7}{0}   % 单位中的中点·按西文处理，避免 kJ·mol 断行
\xeCJKsetcharclass{"2013}{"2013}{0}   % en dash 用西文字体
\setlength{\parskip}{0.15em}

% ---------- 常用宏包 ----------
\usepackage[version=4]{mhchem}
\usepackage{chemfig}
\usepackage{tikz}
\usetikzlibrary{arrows.meta,calc,positioning,shapes.geometric,decorations.pathreplacing,patterns,3d,shadows}
\usepackage{pgfplots}
\pgfplotsset{compat=1.18}
\usepackage{booktabs,tabularx,multirow,array,makecell,colortbl}
\usepackage{enumitem}
\usepackage{multicol}
\usepackage[skins,breakable,most]{tcolorbox}
\usepackage{tasks}
\usepackage{fancyhdr}
\usepackage{xcolor}
\usepackage{caption}
\usepackage{wrapfig}
\usepackage{siunitx}
\sisetup{inter-unit-product=\ensuremath{\cdot},per-mode=symbol}
\usepackage[hidelinks,bookmarksnumbered]{hyperref}

% ---------- 颜色 ----------
\definecolor{mainblue}{RGB}{0,84,147}
\definecolor{lightblue}{RGB}{232,242,251}
\definecolor{deepgreen}{RGB}{20,110,70}
\definecolor{lightgreen}{RGB}{233,246,238}
\definecolor{exorange}{RGB}{196,90,10}
\definecolor{lightorange}{RGB}{254,244,232}
\definecolor{warnred}{RGB}{180,30,40}
\definecolor{lightred}{RGB}{253,236,236}
\definecolor{gkpurple}{RGB}{100,45,140}
\definecolor{lightpurple}{RGB}{244,238,250}
\definecolor{graytext}{RGB}{90,90,90}

% ---------- 章节样式 ----------
\ctexset{
  chapter={
    format={\raggedright\heiti},
    nameformat={\color{mainblue}\Large},
    titleformat={\color{mainblue}\LARGE\bfseries},
    number={\arabic{chapter}},
    name={第,章},
    aftername={\quad},
    beforeskip={0pt},
    afterskip={1.0em},
    pagestyle=fancy
  },
  section={
    format={\Large\heiti\color{mainblue}\raggedright},
    number={\arabic{chapter}.\arabic{section}},
    beforeskip={1.0ex plus .3ex},
    afterskip={0.6ex}
  },
  subsection={
    format={\large\heiti\color{deepgreen}\raggedright},
    number={\arabic{chapter}.\arabic{section}.\arabic{subsection}},
    beforeskip={0.8ex plus .2ex},
    afterskip={0.4ex}
  },
  subsubsection={
    format={\normalsize\heiti\raggedright},
    number={},
    beforeskip={0.5ex},
    afterskip={0.2ex}
  },
  contentsname={目\quad 录}
}
\setcounter{secnumdepth}{2}
\setcounter{tocdepth}{1}

% ---------- 页眉页脚 ----------
\pagestyle{fancy}
\fancyhf{}
\fancyhead[L]{\small\heiti\color{mainblue}\leftmark}
\fancyhead[R]{\small\kaishu\color{graytext}高中化学系统讲义}
\fancyfoot[C]{\small\thepage}
\renewcommand{\headrulewidth}{0.5pt}
\renewcommand{\chaptermark}[1]{\markboth{第\,\thechapter\,章\quad #1}{}}
\fancypagestyle{plain}{\fancyhf{}\fancyfoot[C]{\small\thepage}\renewcommand{\headrulewidth}{0pt}}

% ---------- 列表 ----------
\setlist{nosep,leftmargin=2em}
\setlist[itemize]{label=\textcolor{mainblue}{\small$\blacktriangleright$}}
\setlist[enumerate,1]{label=(\arabic*)}
\settasks{label=\Alph*．,label-width=2.4em,item-indent=2.8em,before-skip=0.1em,after-skip=0.1em,column-sep=0.5em}
\setlength{\emergencystretch}{2.5em}

% ---------- 栏目环境 ----------
\tcbset{
  basestyle/.style={enhanced,breakable,boxrule=0.6pt,arc=1.5mm,left=2mm,right=2mm,top=1mm,bottom=1mm,
    before skip=0.55em,after skip=0.55em,fonttitle=\heiti\bfseries\small}
}
% 概念：核心概念与定义
\newtcolorbox{gainian}[1][]{basestyle,colback=lightblue,colframe=mainblue,
  title={\faicon 核心概念\ifx\relax#1\relax\else：#1\fi},attach boxed title to top left={xshift=3mm,yshift=-2mm},
  boxed title style={colback=mainblue,arc=1mm,boxrule=0pt},top=2.5mm}
% 本质：追本溯源（为什么）
\newtcolorbox{benzhi}[1][]{basestyle,colback=lightgreen,colframe=deepgreen,
  title={追本溯源\ifx\relax#1\relax\else：#1\fi},attach boxed title to top left={xshift=3mm,yshift=-2mm},
  boxed title style={colback=deepgreen,arc=1mm,boxrule=0pt},top=2.5mm}
% 方法：方法提炼
\newtcolorbox{fangfa}[1][]{basestyle,colback=lightorange,colframe=exorange,
  title={方法提炼\ifx\relax#1\relax\else：#1\fi},attach boxed title to top left={xshift=3mm,yshift=-2mm},
  boxed title style={colback=exorange,arc=1mm,boxrule=0pt},top=2.5mm}
% 易错：易错警示
\newtcolorbox{yicuo}[1][]{basestyle,colback=lightred,colframe=warnred,boxrule=0pt,leftrule=2.5pt,arc=0mm,
  title={易错警示\ifx\relax#1\relax\else：#1\fi},coltitle=warnred,colbacktitle=lightred,
  fonttitle=\heiti\bfseries\small,toptitle=0.5mm,bottomtitle=0mm}
% 勘误：原课件勘误
\newtcolorbox{kanwu}[1][]{basestyle,colback=white,colframe=warnred,boxrule=0.6pt,arc=0mm,
  borderline={0.3pt}{1.2pt}{warnred,dashed},
  title={课件勘误\ifx\relax#1\relax\else：#1\fi},coltitle=white,colbacktitle=warnred}
% 真题：真题链接（参数为来源）
\newtcolorbox{zhenti}[1]{basestyle,colback=lightpurple,colframe=gkpurple,
  title={真题链接\quad{\normalfont\small #1}},coltitle=white,colbacktitle=gkpurple}
% 拓展：视野拓展（高考不作要求或仅作了解）
\newtcolorbox{tuozhan}[1][]{basestyle,colback=white,colframe=graytext,boxrule=0.4pt,
  title={视野拓展\ifx\relax#1\relax\else：#1\fi},coltitle=graytext,colbacktitle=white,fonttitle=\heiti\small}
% 导学
\newtcolorbox{daoxue}{basestyle,colback=white,colframe=mainblue,boxrule=1pt,arc=2mm,
  title={本章导学},coltitle=white,colbacktitle=mainblue,fonttitle=\heiti\bfseries}

\def\faicon{}

% 例题：编号随章；上半题目、下半解析
\newcounter{liti}[chapter]
\renewcommand{\theliti}{\thechapter-\arabic{liti}}
\newtcolorbox{liti}[1][]{enhanced,breakable,colback=white,colframe=exorange,boxrule=0.5pt,arc=1mm,
  left=2mm,right=2mm,top=1mm,bottom=1mm,before skip=0.6em,after skip=0.6em,
  segmentation style={dashed,exorange},colbacklower=white,
  before upper={\stepcounter{liti}{\heiti\color{exorange}【例\,\theliti】}\ifx\relax#1\relax\else{\kaishu\small（#1）}\fi\ },
  before lower={\small}}
\newcommand{\silu}{{\heiti\color{exorange}【思路】}}
\newcommand{\jie}{{\heiti\color{exorange}【解析】}}
\newcommand{\daan}{{\heiti\color{exorange}【答案】}}
\newcommand{\xianjing}{{\heiti\color{warnred}【陷阱】}}
\newcommand{\dianbo}{{\heiti\color{deepgreen}【点拨】}}

% 习题与解析
\newlist{xiti}{enumerate}{1}
\setlist[xiti]{label={\textbf{\arabic*}.},leftmargin=1.6em,itemsep=0.35em,topsep=0.3em,resume=xiti}
\newcommand{\zu}[1]{\par\smallskip\noindent{\heiti\color{mainblue}\large #1}\par\smallskip}
\newcommand{\laiyuan}[1]{{\kaishu\small\color{graytext}（#1）}}
\newenvironment{jiexilist}{\small\begin{enumerate}[label={\textbf{\arabic*}.},leftmargin=1.6em,itemsep=0.3em,topsep=0.2em]}{\end{enumerate}}
\newcommand{\ans}[1]{{\heiti 答案：}#1\quad}

% 填空线、括号
\newcommand{\tk}[1][3em]{\underline{\hspace{#1}}}
\newcommand{\xz}{（\hspace{2em}）}

% 化学常用
\newcommand{\kjmol}{\,kJ·mol$^{-1}$}
\newcommand{\molL}{\,mol·L$^{-1}$}
\newcommand{\dH}{\Delta H}
\newcommand{\ksp}{K_{\mathrm{sp}}}
\newcommand{\ka}{K_{\mathrm{a}}}
\newcommand{\kb}{K_{\mathrm{b}}}
\newcommand{\kw}{K_{\mathrm{w}}}
\newcommand{\kh}{K_{\mathrm{h}}}
\newcommand{\kp}{K_{p}}
\newcommand{\NA}{N_{\mathrm{A}}}
\newcommand{\vz}{v_{\text{正}}}
\newcommand{\vn}{v_{\text{逆}}}
\newcommand{\kz}{k_{\text{正}}}
\newcommand{\kn}{k_{\text{逆}}}
\newcommand{\eqm}{\rightleftharpoons}

% 小节标题“章末”样式
\newcommand{\zhangmo}[1]{\section*{#1}\addcontentsline{toc}{section}{#1}}

% 表格行距
\renewcommand{\arraystretch}{1.25}
\setlength{\tabcolsep}{4pt}

% chemfig 设置
\setchemfig{atom sep=1.8em,bond offset=1pt,double bond sep=2.2pt,angle increment=30}

% 能量-历程图快捷命令：\energyprofile{起点}{峰}{终点}
\newcommand{\axisEP}[2]{%
  \draw[-Stealth] (0,0)--(#1,0) node[below left]{\small 反应进程};
  \draw[-Stealth] (0,0)--(0,#2) node[left]{\small 能量};}

% 热化学方程式中的焓变：\dHv{-285.8}
\newcommand{\dHv}[1]{\quad\ensuremath{\Delta H=#1\ \mathrm{kJ\cdot mol^{-1}}}}
% 题干中的选项（行内四选项，自动排版）
\newcommand{\xx}[4]{\begin{tasks}(4)\task #1 \task #2 \task #3 \task #4\end{tasks}}
\newcommand{\xxb}[4]{\begin{tasks}(2)\task #1 \task #2 \task #3 \task #4\end{tasks}}
\newcommand{\xxa}[4]{\begin{tasks}(1)\task #1 \task #2 \task #3 \task #4\end{tasks}}
